\chapter{Lecture 9: Classes, Constructors and Member Initialization Lists}
\begin{center}
\textit{February 1, 2018}
\end{center}

\section{Separate Compilation: Rules}
Last time we fixed the ``redefinition of struct Vec'' error
(\code{lectures/separate/example3}) with an include guard
(\code{lectures/separate/example4}). The rules:
\begin{enumerate}
  \item All \code{.h} files should have an include guard.
  \item Never compile \code{.h} files.
  \item Never \code{\#include} \code{.cc} files.
\end{enumerate}

\begin{important}[No \code{using namespace std} in headers]
Never put \code{using namespace std;} in a header file: every file that includes
your header would be forced to have it too (possibly causing name clashes). In a
header, refer to \code{std::string}, \code{std::istream}, \code{std::ostream},
etc.\ by their full names.
\end{important}

\section{C++ Classes}
The big OO innovation: we can put functions inside a struct.

\begin{lstlisting}[style=cpp]
// student.h
#ifndef STUDENT_H
#define STUDENT_H
struct Student {
    int assns, mt, final;
    float grade();
};
#endif

// student.cc
#include "student.h"
float Student::grade() {
    return assns * 0.4 + mt * 0.2 + final * 0.4;
}

// main.cc
Student billy{80, 50, 75};
cout << billy.grade() << endl;   // 70
\end{lstlisting}

\code{::} is the \textbf{scope resolution operator}. \code{Student::grade} means
``the function \code{grade} in the scope of the class \code{Student}''.

\begin{definition}
\begin{itemize}
  \item A \textbf{class} is (essentially) a struct type that can contain
        functions. In C++ every struct is a class.
  \item An \textbf{object} is an instance of a class: a value of type
        \code{Student} is a \code{Student} object; \code{billy} is an object.
  \item A function inside a class is a \textbf{member function} or
        \textbf{method}.
\end{itemize}
\end{definition}

A method can only be called on an object of its class, and has access to the
fields of the object on which it was called. In \code{grade}, \code{assns}
means \code{billy.assns} when we call \code{billy.grade()}.

\subsection{The \code{this} Pointer}
Every method has a hidden parameter called \code{this}, which points to the
object on which the method was called. In \code{billy.grade()},
\code{this == \&billy}. So the method above is really
\begin{lstlisting}[style=cpp]
float Student::grade() {
    return this->assns * 0.4 + this->mt * 0.2 + this->final * 0.4;
}
\end{lstlisting}
\begin{quote}
\textit{``I like to use \code{this} when I need to use \code{this}.''} --- Nomair
Naeem
\end{quote}
(\code{this} is a constant pointer: you cannot re-point it.)

\section{Initializing Objects}
\code{Student billy\{80, 50, 75\};} uses C-style field-by-field initialization.
This is limited: we cannot validate the values, compute anything, or allocate
memory. In C++, you can write special methods called \textbf{constructors}
(ctors) to initialize objects.

\begin{lstlisting}[style=cpp]
struct Student {
    int assns, mt, final;
    Student(int assns, int mt, int final);
    float grade();
};

Student::Student(int assns, int mt, int final) {
    this->assns = assns < 0 ? 0 : assns;   // we can sanity-check now
    this->mt = mt < 0 ? 0 : mt;
    this->final = final < 0 ? 0 : final;
}

Student billy{80, 50, 75};   // calls the constructor
\end{lstlisting}
A constructor has the same name as the class and no return type. Here the
parameter names shadow the field names, so we need \code{this->} to reach the
fields.

\begin{itemize}
  \item \code{Student billy\{80, 50, 75\};} --- stack object, calls the ctor
        with the three arguments.
  \item \code{Student *pBilly = new Student\{80, 50, 75\};} --- heap object, same
        ctor.
  \item \code{Student billy = Student\{80, 50, 75\};} also works.
\end{itemize}

The \code{\{...\}} syntax is \textbf{uniform initialization} (C++11), and can be
used for all types: \code{int x\{5\};}, \code{string s\{"hello"\};}. If a class
has a ctor, \code{\{...\}} calls it; otherwise (for plain structs) it does
field-by-field initialization.

\subsection{Advantages of Constructors}
Constructors are functions, so we get default arguments, overloading and
sanity checks:
\begin{lstlisting}[style=cpp]
struct Student {
    int assns, mt, final;
    Student(int assns = 0, int mt = 0, int final = 0);
};

Student s1{60, 70};   // final = 0
Student s2;           // all zero: calls Student(0, 0, 0)
\end{lstlisting}

\subsection{The Default Constructor}
Every class comes with a \textbf{default constructor} (a ctor with no
arguments) if you don't write any constructor. It default-constructs every
field that is an object, and leaves primitive fields (\code{int}, pointers)
\emph{uninitialized}.

\begin{important}[You lose the free default ctor]
As soon as you write \emph{any} constructor, the compiler no longer provides the
default constructor.
\begin{lstlisting}[style=cpp]
struct Vec {
    int x, y;
    Vec(int x, int y) : x{x}, y{y} {}
};
Vec v;          // error: no default ctor
Vec w{1, 2};    // OK
\end{lstlisting}
\end{important}

\section{Member Initialization Lists}
Consider a class with a constant or reference field:
\begin{lstlisting}[style=cpp]
int z;
struct MyStruct {
    const int myConst;   // must be initialized
    int &myRef;          // must be initialized
};
\end{lstlisting}
We can't give them default values in the class (every object could need a
different value), and we can't assign to them in the ctor body. Why? Because of
how objects are built.

\begin{important}[Steps of object construction]
\begin{enumerate}
  \item Space is allocated (stack or heap).
  \item Fields are initialized: each field that is an object has its
        constructor called (default ctor, unless told otherwise).
  \item The constructor body runs.
\end{enumerate}
\end{important}

By the time the body runs, the fields have already been initialized, so
\code{myConst = 5;} in the body would be an assignment, not an initialization.
To hook into step 2, use the \textbf{member initialization list} (MIL):
\begin{lstlisting}[style=cpp]
struct Student {
    const int id;          // a student's id never changes
    int assns, mt, final;
    Student(int id, int assns, int mt, int final);
};

Student::Student(int id, int assns, int mt, int final)
    : id{id}, assns{assns}, mt{mt}, final{final} {}
\end{lstlisting}
In \code{id\{id\}}, the name outside the braces is always the \emph{field}, and
the expression inside the braces is looked up normally (so here it is the
parameter). No \code{this->} needed.

Facts about the MIL:
\begin{itemize}
  \item You \emph{must} use it to initialize \code{const} fields, reference
        fields, and object fields with no default constructor.
  \item It can be used for all fields.
  \item Fields are initialized in the order they are \emph{declared in the
        class}, regardless of the order in the MIL (compile with \code{-Wall}
        to be warned about mismatched orders).
  \item It can be more efficient than assigning in the body: for an object
        field like a \code{string}, assigning in the body means it is first
        default-constructed and then assigned (two steps), whereas the MIL
        constructs it directly with the right value (one step).
\end{itemize}

You can also give fields default values in the class (C++11); the MIL overrides
them:
\begin{lstlisting}[style=cpp]
struct Vec {
    int x = 0, y = 0;
    Vec() {}                          // x = y = 0
    Vec(int x) : x{x} {}              // y = 0
    Vec(int x, int y) : x{x}, y{y} {}
};
\end{lstlisting}

\section*{Summary}
\begin{itemize}
  \item Include guards in headers, never compile headers, never include
        \code{.cc} files, no \code{using namespace std} in headers.
  \item Classes are structs with methods; methods get a hidden \code{this}
        pointer; \code{Class::method} defines a method outside the class.
  \item Constructors initialize objects; writing any ctor removes the free
        default ctor.
  \item Construction: allocate, initialize fields (MIL), run body. Use the MIL
        for \code{const}/reference fields and fields without default ctors;
        fields are initialized in declaration order.
\end{itemize}
